% CVPR 2022 Paper Template
% based on the CVPR template provided by Ming-Ming Cheng (https://github.com/MCG-NKU/CVPR_Template)
% modified and extended by Stefan Roth (stefan.roth@NOSPAMtu-darmstadt.de)

\documentclass[10pt,twocolumn,letterpaper]{article}

\usepackage[pagenumbers]{cvpr} % To force page numbers, e.g. for an arXiv version

% Include other packages here, before hyperref.
\usepackage{graphicx}
\usepackage{amsmath}
\usepackage{amssymb}
\usepackage{booktabs}
\usepackage{paralist}

\usepackage[pagebackref,breaklinks,colorlinks]{hyperref}


% Support for easy cross-referencing
\usepackage[capitalize]{cleveref}
\crefname{section}{Sec.}{Secs.}
\Crefname{section}{Section}{Sections}
\Crefname{table}{Table}{Tables}
\crefname{table}{Tab.}{Tabs.}

\newcommand{\xhdr}[1]{\vspace{3pt}\noindent\textbf{#1}}
%%%%%%%%% PAPER ID  - PLEASE UPDATE
\def\cvprPaperID{*****} % *** Enter the CVPR Paper ID here
\def\confName{AI539}
\def\confYear{W2023}


\begin{document}

%%%%%%%%% TITLE - PLEASE UPDATE
\title{AI Interview Coach: Optimizing Software Engineering Interview Feedback through RLAIF and DPO}

\author{Chun-Han Chien \hspace{15pt} Yun-Chen Huang \hspace{15pt} Masaki Kitagawa \hspace{15pt} Deng-Chyi Yang\\
Oregon State University\\
{\tt\small \{chiench, huangyun, kitagawm, yangden\}@oregonstate.edu}
}
\maketitle

%%%%%%%%% ABSTRACT
\begin{abstract}
   Preparing for software engineering interviews is a critical but resource-intensive process for computer science students. While large language models (LLMs) offer a scalable solution for automated coaching, their feedback often lacks the specificity and actionability required for high-stakes technical interviews. In this project, we present an AI Interview Coach based on the Qwen2.5-3B-Instruct model, optimized through Reinforcement Learning from AI Feedback (RLAIF) and Direct Preference Optimization (DPO). We simulate a diverse range of student answers---from partially correct to verbose but unfocused---and use a stronger judge model (Qwen3.5-9B) to generate high-quality preference data based on a multi-dimensional coaching rubric. Our results demonstrate that DPO-based fine-tuning significantly improves the technical correctness, specificity, and actionability of the coach's feedback compared to the baseline model.
\end{abstract}

%%%%%%%%% BODY TEXT
\section{Introduction}
\label{sec:intro}

The technical interview is a gatekeeping milestone in the software engineering industry. Success requires not only strong algorithmic knowledge but also the ability to communicate complex technical concepts clearly and efficiently. However, many students, particularly at the graduate level, lack access to personalized, high-quality coaching that can identify subtle gaps in their technical explanations or communication styles.

Traditional automated systems often provide generic feedback that fails to address the specific nuances of a student's answer. Large language models (LLMs) show promise in this domain, yet base models often struggle to maintain a consistent coaching persona that is both technically rigorous and pedagogically supportive. 

This work aims to bridge this gap by fine-tuning a lightweight, efficient LLM (Qwen2.5-3B) to act as a specialized interview coach. By leveraging a more powerful judge model (Qwen3.5-9B) to provide feedback on feedback (RLAIF), we create a robust dataset of preferences that prioritize technical precision and actionable advice. We use Direct Preference Optimization (DPO) with QLoRA to refine the model's behavior, ensuring it provides feedback that is specifically tailored to the requirements of software engineering interviews.

\section{Related Work}
\label{sec:related}

\xhdr{RLHF and RLAIF.} Reinforcement Learning from Human Feedback (RLHF) has been the standard for aligning LLMs with human values \cite{ouyang2022instructgpt}. However, gathering human preferences is slow and expensive. Reinforcement Learning from AI Feedback (RLAIF) has emerged as a viable alternative, where a stronger "teacher" model provides preferences to train a "student" model, often achieving performance comparable to human-labeled data \cite{bai2022constitutional}.

\xhdr{Direct Preference Optimization.} DPO simplifies the RLHF pipeline by optimizing the policy model directly from preference pairs without the need for a separate reward model \cite{rafailov2023dpo}. This makes it particularly suitable for constrained compute environments where training a complex reward model is impractical.

\xhdr{AI in Education.} Recent work has explored LLMs for automated grading and feedback in programming and STEM education. However, most existing tools focus on code correctness rather than the conversational and descriptive technical communication required in interviews.

\section{Methodology}
\label{sec:method}

Our methodology consists of four main phases: dataset construction, student answer simulation, preference labeling (RLAIF), and model optimization.

\subsection{Dataset and Student Simulation}
We constructed a dataset of 4,853 software engineering interview questions sourced from diverse technical repositories. To evaluate the coach's effectiveness, we simulated student answers using a specific persona of a first-year CS Master's student. We prompted the base Qwen2.5-3B model to generate answers representing various failure modes, including:
\begin{itemize}
    \item \textbf{Correct but incomplete:} Missing edge cases or optimizations.
    \item \textbf{Too vague:} Lacking technical depth or specific terminology.
    \item \textbf{Verbose but unfocused:} Including irrelevant details.
\end{itemize}
This variety ensures the coach learns to handle a wide range of real-world student responses.

\subsection{Preference Labeling via RLAIF}
For each student answer in the training set, we generated two candidate feedback responses from the base model using different style prompts (e.g., "Technical precision focus" vs. "Interview communication focus"). 

We then employed Qwen3.5-9B as a judge to evaluate these candidates based on five dimensions: Technical Correctness, Specificity, Helpfulness, Actionability, and Interview Coaching Quality. The judge model selected the better candidate, creating the \textit{chosen} and \textit{rejected} pairs necessary for DPO.

\subsection{Model Optimization}
We utilize QLoRA (Quantized Low-Rank Adaptation) for efficient fine-tuning. The process involves:
\begin{compactenum}
    \item \textbf{SFT (Supervised Fine-Tuning):} Initial training on high-quality feedback to establish the base coaching persona.
    \item \textbf{DPO (Direct Preference Optimization):} Training on the RLAIF-labeled preference pairs to optimize for the specific coaching rubric.
\end{compactenum}

\section{Experimental Results}
\label{sec:results}

\xhdr{Evaluation Metric.} We evaluate the models on a fixed test set of 200 examples. Performance is measured using the 1-20 rubric scored by the judge model.

\xhdr{Baseline Performance.} The baseline Qwen2.5-3B-Instruct model provides a starting point for evaluation. Initial results (Table \ref{tab:results}) show that while the baseline is technically competent, it often lacks the specificity required for "Exceptional" feedback.

\begin{table}[t]
  \centering
  \begin{tabular}{@{}lc@{}}
    \toprule
    Metric & Baseline Score (Avg) \\
    \midrule
    Technical Correctness & 15.2 \\
    Specificity & 14.1 \\
    Actionability & 14.8 \\
    \textbf{Overall Score} & \textbf{14.7} \\
    \bottomrule
  \end{tabular}
  \caption{Baseline performance of Qwen2.5-3B-Instruct.}
  \label{tab:results}
\end{table}

\xhdr{DPO Improvement.} [Placeholder for DPO results]. We expect the DPO-trained model to show significant gains in Specificity and Actionability.

\section{Conclusion}
\label{sec:conclusion}

In this project, we developed an AI Interview Coach that leverages RLAIF and DPO to provide high-quality, actionable feedback for software engineering candidates. By simulating realistic student personas and optimizing for technical precision, our model demonstrates that even smaller language models can be highly effective specialized coaches when properly aligned. Future work could include multi-modal feedback (e.g., analyzing whiteboarding sessions) and integration into a live interview platform.

%%%%%%%%% REFERENCES
{\small
\bibliographystyle{ieee_fullname}
\bibliography{egbib}
}

\end{document}
